\section{Normalized global propagation on the actual cubic chart}
\label{sec:cubic-normalized-global-propagation}

The local character conditions in this section are the normalized
conditions of AC-G101 and AC-G104.  Constant nontrivial cubic ratios
are excluded.  The original one-character \texttt{Nice} predicate is
not substituted for this predicate.  We prove the global geometry and
its local-to-global implication, retaining the precise arithmetic
seed and rank-two support inputs when they are used.

Use the actual rings $B\subset A\subset R$ of AC-G94--AC-G95 and
put $X=\operatorname{Spec}A$.  To make the triangular convention
unambiguous, write
\[
 U=(r_{ij}(g):i<j),\qquad L=(r_{ij}(g):i>j),\qquad
 E=(r_{ij}(g):i-j\text{ odd}).
 \tag{G108.1}
\]
Thus $U=T_-$ and $L=T_+$ in AC-G101.  Define
\[
 \begin{split}
 X_{\mathrm{off}}&=X\setminus\bigl(V(U)\cup V(L)\bigr),\\
 X_{\mathrm I}&=X_{\mathrm{off}}\setminus V(E),\qquad
 X_{\mathrm E}=X_{\mathrm{off}}\cap V(E),\\
 \mathcal C&=V(\mathfrak m_BA)\setminus V(L),\qquad
 \mathcal C_{\mathrm{cl}}=\{x\in\mathcal C:x\text{ is closed in }X\}.
 \end{split}
 \tag{G108.2}
\]
The set $\mathcal C$ is exactly the punctured lower locus proved
connected in AC-G95-lower-locus, regardless of the names of the two
triangular ideals there.  Its points kill $U$.
For an ideal $J\subset B$, write
\[
 \operatorname{Fin}_{\theta}(J):\quad
 \sqrt{J+\lambda_{\theta}}=\mathfrak m_B,
 \qquad \operatorname{Good}_J=V(JA).
 \tag{G108.3}
\]
Here $\lambda_{\theta}$ is the joint-residual-kernel weight ideal
of AC-G101, with the coefficient uniformizer included.  This is the
actual finite-weight condition proved for the native ideal by AC-G102.

\begin{theorem}[AC-G108-geometry: component access and seed localization]
\label{AC-G108-geometry}
Every irreducible component of $X$ meeting $X_{\mathrm I}$ meets
$\mathcal C$, and does so at a closed point of $X$.  Every nonempty
locally closed subset of $\mathcal C$ contains a point of
$\mathcal C_{\mathrm{cl}}$.  Under the actual filled adjacent and
central class hypotheses of AC-G93-fill,
\[
 \mathcal C\text{ is connected},\qquad
 \mathcal C\subset
 \overline{\mathcal C_{\mathrm{cl}}\setminus V(E)}.
 \tag{G108.4}
\]
Suppose in addition that there is a prime $\xi\in X_{\mathrm I}$
with $Z=\xi\cap B$.  Then there is a closed $x\in\mathcal C$
and a prime $P$ of its actual completion $C_x$ such that
\[
 P\cap B=Z,\quad P\in X_{\mathrm I}(C_x),\quad
 \dim C_x/P=\dim B/Z+1.
 \tag{G108.5}
\]
Here $X_{\mathrm I}(C_x)$ is defined by the extensions of the three
ideals in (G108.1).  In particular the actual strong seed
$\dim B/Z=3d+1$, $J\subset Z$ gives a good Type-I local prime of
dimension $3d+2$.  This assertion does not deduce such a seed from
one characteristic-zero point or from an arbitrary quotient of a
weight-torsionfree ring.
\end{theorem}

\begin{proof}
\emph{1. Finite homogeneous pieces.}
Every $A_j$ is a finite $B$-module.  Indeed AC-G95 writes $A$ as a
$B$-algebra on finitely many homogeneous coordinates of weights
$0,\pm1,\pm2$, and $A_0=B$.  Remove zero-weight factors from a
coordinate monomial.  The four removable patterns are
$(1,-1),(2,-2),(1,1,-2),(2,-1,-1)$.  A remaining monomial of
specified weight $j$ has bounded degree: with one sign its degree
is at most $|j|$; with both signs the only possible pairs are
$\{1,-2\}$ and $\{2,-1\}$, and the lighter opposite-sign factor
has exponent at most one.  The weight equation bounds the other
exponent.  There are only finitely many remainder monomials.
The removed factors lie in $B$, proving finiteness.
For a homogeneous ideal $\eta$, the degree-$j$ part of $A/\eta$
is consequently a finite $B$-module as well.

\emph{2. Homogeneous primes meet the punctured closed fiber.}
Let $\eta$ be homogeneous and prime, with $L\not\subset\eta$.
Set $D=A/\eta$.  Suppose every negative-weight coordinate becomes
nilpotent in $D/\mathfrak m_BD$.  There are finitely many such
coordinates.  Their exponents in a nonzero coordinate monomial
in this quotient are then bounded.  Positive-weight coordinates
cannot lower the degree.  Hence for some $m$,
\[
 (D/\mathfrak m_BD)_j=0\quad(j<-m).
 \tag{G108.6}
\]
Since $\mathfrak m_B$ has degree zero, this quotient piece is
$D_j/\mathfrak m_BD_j$.  Finiteness from Step 1 and Nakayama give
$D_j=0$ for $j<-m$.  But some strictly lower entry $a$ survives
in the domain $D$ and has negative degree.  Every power of $a$
is nonzero, contrary to (G108.6).  Thus $L$ is not contained in
$\sqrt{\eta+\mathfrak m_BA}$.
There is a prime over $\eta+\mathfrak m_BA$ at which $L$ survives.
The quotient $A/\mathfrak m_BA$ is a finite-type $k$-algebra and
is Jacobson.  In the nonempty open subset of
$V(\eta+\mathfrak m_BA)$ where this lower entry survives, choose
a closed point.  It is maximal in $A$ and belongs to $\mathcal C$.

Every minimal prime of $A$ is homogeneous.  For a prime $\eta$,
map $A$ to $\operatorname{Frac}(A/\eta)[t,t^{-1}]$ by
$a_j\mapsto\bar a_jt^j$ on homogeneous elements.  Its kernel
is the homogeneous core of $\eta$, is prime, and is contained in
$\eta$.  Minimality forces equality.  If an irreducible component
meets $X_{\mathrm I}$, its minimal prime does not contain $L$.
Step 2 therefore proves the component-access assertion.

\emph{3. Closed points and density.}
The set $\mathcal C$ is open in the closed finite-type $k$-fiber.
A nonempty locally closed subset of that fiber contains a closed
point, by the Jacobson property just used.  Its residue field is
finite over $k$, so its inverse image in $A$ is a maximal ideal.
This proves the stated closed-point assertion.
AC-G95-lower-locus proves connectedness and
$\mathcal C\subset\overline{\mathcal C\setminus V(E)}$ by actual
polynomial families.  Every open neighborhood of a point in
$\mathcal C\setminus V(E)$ contains a closed point in that same
set, by the locally closed assertion.  Taking closures twice
proves (G108.4).  No irreducibility of the lower mixed-cup locus
is required.

\emph{4. The strong seed.}
Apply the homogeneous map of Step 2 to $\xi$, and call its kernel
$\eta$.  Then $\eta\cap B=Z$.  For a homogeneous element $a$,
$a\in\eta$ if and only if $a\in\xi$.  Every entry generating
$U,L,E$ is homogeneous, so $\eta\in X_{\mathrm I}$.
The graded dimension argument in AC-G101 gives
\[
 \operatorname{trdeg}_{\operatorname{Frac}(B/Z)}
            \operatorname{Frac}(A/\eta)\leq1.
 \tag{G108.7}
\]
A nonzero homogeneous element of nonzero degree is transcendental
over $\operatorname{Frac}(B/Z)$: after clearing coefficients,
any polynomial relation has terms in distinct degrees, each of
which must vanish; the domain condition makes every coefficient
zero.  Such an element exists because $\eta$ is off triangular.
Consequently the transcendence degree in (G108.7) is exactly one.
Step 2 gives a closed $x\in\mathcal C$ containing $\eta$.

We use the precise generic commutative-algebra statement
\texttt{LocDim.lean}, private theorem \texttt{dim\_gen}, second
conclusion.  Its entire containing proof is retained in the source
receipt.  Its hypotheses are: a complete local quotient $B$ over
the complete coefficient DVR, a finite-type algebra $A$, a finite
flat coefficient extension, a closed point over $\mathfrak m_B$,
and completion at a point over it.  These are exactly AC-G95 and
the actual construction of $C_x$ in AC-G101.  Its conclusion
lifts $\eta$ to a prime $P\subset C_x$ with the same contraction
and dimension $\dim B/Z+t$, where $t$ is (G108.7).
The proof uses the dimension formula and equidimensionality of
these completions, not a dimension bound for arbitrary formal
fibers.  Here $t=1$.  Contraction preserves noncontainment of all
three extended ideals.  This proves (G108.5) and the final assertion.
\end{proof}

The following theorem states all remaining arithmetic inputs as
actual ideal and prime conditions.  It has no implicit reference
to an old residual tuple.

\begin{theorem}[AC-G108-propagation: normalized local-to-global theorem]
\label{AC-G108-propagation}
Assume the actual hypotheses of AC-G93--AC-G105, including
\[
 h+a+z+1<d,\qquad d_v>2\quad(v\mid p),
 \tag{G108.8}
\]
and retain the native N4 and prime-core premises used in
AC-G104-prime.  Let $J,J_{\mathrm E}\subset B$ satisfy
$\operatorname{Fin}_{\theta}(J)$ and
$\operatorname{Fin}_{\theta}(J_{\mathrm E})$.  Assume
\[
 \begin{array}{ll}
 \text{(E)}&J_{\mathrm E}\subset\mathfrak p\cap B
                   \quad(\mathfrak p\in X_{\mathrm E});\\
 \text{(S)}&\text{there is }\xi\in X_{\mathrm I},\quad
        Z=\xi\cap B,\quad J\subset Z,\quad\dim B/Z=3d+1.
 \end{array}
 \tag{G108.9}
\]
Here the native prime-core premises mean precisely the two
normalized conclusions (G104.8) and its finite-chain analogue
for the ideal $J$: normalized nice Type-I primes propagate
$J$ downwards, and finite chains of normalized nice endpoint
primes sharing thick primes propagate $J$ between Type-I primes.
They may be obtained from the actual native ideal by AC-G104;
they are not asserted for an arbitrary ideal merely satisfying
(G108.3).
Then
\[
 X_{\mathrm I}\subset\operatorname{Good}_J.
 \tag{G108.10}
\]
Moreover $J_{\mathrm{off}}=J\cap J_{\mathrm E}$ satisfies
$\operatorname{Fin}_{\theta}(J_{\mathrm{off}})$ and
\[
 J_{\mathrm{off}}\subset\mathfrak p\cap B
                 \quad(\mathfrak p\in X_{\mathrm{off}}).
 \tag{G108.11}
\]
Thus the global geometric obligation is discharged under the
explicit strong seed and the uniform even/off-triangular ideal.
A finite ideal for one endpoint branch does not suffice for (E).
The arithmetic production of (E) and (S) is separate from this theorem.
\end{theorem}

\begin{proof}
Put $D=3d+1$.  At a closed $x\in\mathcal C$, let $C=C_x$.
All ideals and Type-I/endpoint predicates in the following local
argument are extended from (G108.1).  In particular $LC=C$, since
some lower entry is outside $x$.

\emph{1. Normalized curve avoidance.}
If $Q\subset C$ is prime, $JC\subset Q$ and $\dim C/Q\geq D$,
there is a normalized nice prime $q\supset Q$.  To verify this
with the changed tuple, use the finite family of bad ideals
\[
 UC,\qquad JC+(\varpi)+\widehat c_{v,ij}\quad(i\ne j).
 \tag{G108.12}
\]
AC-G98 and (G108.8) bound the first quotient dimension strictly
below $3d$.  AC-G101, using $\operatorname{Fin}_{\theta}(J)$,
bounds each other quotient dimension by
$\max\{2d+h+a+z+1,3d+2-d_v\}<3d$.
The exact tuple-free theorem \texttt{NicePrimeEx.nicePrimeEx}
applies with $N=3d$, $\varpi\in\mathfrak m_C$, and
$\dim C/Q\geq N+1$.  It supplies a prime $q\supset Q$ with
$\dim C/q=1$, $\varpi\in q$, and none of (G108.12) contained
in $q$.  Since $JC$ and $\varpi$ already lie in $q$, no
$\widehat c_{v,ij}$ is contained in $q$.  Also $UC\not\subset q$
and $LC=C$.  This is exactly the normalized nice definition.
The coefficient prime is in $q$ because it is a unit times a
positive power of $\varpi$.
The same assertion holds with $J_{\mathrm E}$ replacing $J$.
We have checked the applied parameters of the source prime-avoidance
lemma; no old \texttt{Frame.Big} or raw collision ideal is imported.

If a thick prime $e$ lies in the actual Sharp component set and
$Q\supset e$ has $\dim C/Q\geq D$, (E) applies to $e\cap A$.
Indeed contraction preserves its endpoint predicate.  Therefore
$J_{\mathrm E}C\subset e\subset Q$.  The preceding construction
with $J_{\mathrm E}$ gives a normalized nice $q\supset Q$.
As $EC\subset e$, this $q$ is endpoint type.  This explains the
need for (E) on every even/off-triangular prime, without assuming
in advance that a thick component is good for $J$.

\emph{2. Propagation among local components.}
AC-G105 gives an ideal $s\subset C$.  Every prime outside $V(EC)$
contains $s$.  Every minimal prime over $s$ has dimension at least
$D+1$, and any two are joined by a finite chain whose consecutive
intersection quotients have dimension at least $D$.  Such a
minimal prime is Type I or thick.  Every Type-I minimal prime
of $C$ is a minimal prime over $s$.
The unit-ideal case has no such primes; the assertions below
are then vacuous unless an actual Type-I prime is supplied.

Call a local prime good if it contains $JC$.  Goodness goes up
under inclusion, and the Type-I property goes down.  At any
normalized nice prime $q$, the following implication follows
from the retained normalized prime-core conclusions:
\[
 \begin{gathered}
 \text{some good Type-I prime lies below }q\quad\Longrightarrow\\
 \text{every Type-I prime below }q\text{ is good}.
 \end{gathered}
 \tag{G108.13}
\]
If $q$ is Type I, it is itself good and the first conclusion
of AC-G104 applies.  If $q$ is endpoint type, apply the second
conclusion to the length-zero chain.  Every normalized nice
prime is one of these two types, by whether it contains $EC$.

Suppose $I$ is good, Type I, and $\dim C/I\geq D$.
Choose a minimal prime $P_0\subset I$.  It is Type I and a
minimal prime over $s$.  Step 1 gives a normalized nice
$q\supset I$, and (G108.13) makes $P_0$ good.
We propagate along the finite component chains from $P_0$.
Use the following invariant on a minimal prime $P$ over $s$:
either $P$ is good Type I; or $P$ is thick and there is a finite
chain of normalized nice endpoint primes beginning at $q_0$
and ending at $q_1\supset P$, with a good Type-I prime below
$q_0$, such that successive endpoint primes share a thick
minimal prime over $s$.
The invariant holds at $P_0$.

Here is its passage across one adjacent pair $P,P'$.
The dimension bound on $C/(P+P')$ gives a prime $Q\supset P+P'$
with $\dim C/Q\geq D$; choose a maximal-dimensional irreducible
component of that quotient.  If $P$ is good Type I, $Q$ is good,
so Step 1 gives nice $q\supset Q$, and (G108.13) makes every
Type-I prime below $q$ good.  Thus $P'$ is good if Type I.
If $P'$ is thick, use the second paragraph of Step 1 on $P'$
and $Q$ to choose a nice endpoint $q'\supset Q$; the good
Type-I prime $P$ lies below $q'$.  The length-zero endpoint
chain at $q'$ establishes the invariant at $P'$.

If $P$ is thick with its invariant chain ending at $q_1$, use
Step 1 on $P$ and $Q$ to obtain a nice endpoint $q\supset Q$.
Append the link from $q_1$ to $q$ through $P$.
The second prime-core conclusion makes every Type-I prime
below $q$ good, so it handles Type-I $P'$.
If $P'$ is thick, choose a nice endpoint $q'\supset Q$ using
$P'$ and Step 1.  Append instead the link from $q_1$ to $q'$
through $P$.  This endpoint contains $P'$ and gives its invariant.
These are all cases.  Finite induction along the component chain
therefore makes every Type-I minimal prime of $C$ good.
We have proved the actual local conclusion
\[
 \begin{gathered}
 I\text{ good Type I},\quad\dim C/I\geq D\quad\Longrightarrow\\
 \text{all Type-I minimal primes of }C\text{ are good}.
 \end{gathered}
 \tag{G108.14}
\]
This spells out the tuple-free \texttt{LocProp.locProp} argument
with the verified normalized avoidance input.

\emph{3. From completions to components.}
The actual map $A\to C_x$ is flat: it is a finite flat unramified
coefficient extension, localization, and Noetherian local
completion.  Its closed point contracts to $x$.  Minimal primes
of $C_x$ contract to minimal primes of $A$, by going down;
every component of $X$ through $x$ is reached by a minimal prime
of $C_x$, using going down below the closed point and then a
minimal prime.  Equivalently, the closure of the image of each
minimal prime is a component of $X$, and all components through
$x$ occur.  These are exactly the generic statements proved in
\texttt{ChartOfFlat.chartOfFlat}; all hypotheses have been checked.
The extended Type-I predicate is the inverse image of
$X_{\mathrm I}$, and goodness is the inverse image of
$\operatorname{Good}_J$.

Suppose two components $Y,Y'$ meeting $X_{\mathrm I}$ contain
$x\in\mathcal C_{\mathrm{cl}}$, and $Y$ is good.  Choose the
corresponding minimal primes $P,P'\subset C_x$.  Their generic
points are Type I, since $X_{\mathrm I}$ is open and each
component meets it.  The prime $P$ is good and has dimension
at least $3d+2=D+1$ by AC-G101.  Apply (G108.14) to $P$.
Then $P'$ is good, and closedness of $V(JA)$ implies all of $Y'$
is good.  This proves propagation between components meeting
at a closed point of the actual core.

If $x\in\mathcal C_{\mathrm{cl}}\setminus V(E)$, then $EC_x=C_x$
and AC-G105 uses $s=0$.  Its dichotomy says every minimal prime
of $C_x$ is Type I; a thick prime would contain $EC_x$.
Consequently $x$ lies on a component meeting $X_{\mathrm I}$.
Finally (S) and (G108.5) give a good Type-I prime of a core
completion of dimension $3d+2$.  Apply (G108.14), and contract
a minimal prime below it.  Its component meets $X_{\mathrm I}$
and is entirely good.  This constructs a good component; mere
nonemptiness of the native support was not used as a substitute.

\emph{4. The connected core forces all relevant components good.}
There are finitely many irreducible components of $X$.  Let
$Y_{\mathrm g}$ be the union of those meeting $X_{\mathrm I}$
and contained in $V(JA)$; let $Y_{\mathrm b}$ be the union of
those meeting $X_{\mathrm I}$ and not so contained.  Both unions
are closed.  Step 3 and (G108.4) show
$\mathcal C\subset Y_{\mathrm g}\cup Y_{\mathrm b}$.
The good union meets the core by the constructed good component
and AC-G108-geometry.  If there is a bad component, the bad union
also meets the core by the same component-access theorem.
Connectedness of $\mathcal C$ forces
\[
 Y_{\mathrm g}\cap Y_{\mathrm b}\cap\mathcal C\ne\varnothing.
 \tag{G108.15}
\]
The closed-point assertion in AC-G108-geometry gives a closed
$x$ in this intersection.  At $x$ a good relevant component
and a bad relevant component meet, contradicting the propagation
proved in Step 3.  Thus there are no bad components.
Every Type-I point lies on a component meeting the Type-I open
set, so (G108.10) follows.

Every off-triangular point is Type I or endpoint type.
For the former (G108.10) puts $J$ in its contraction; for the
latter (E) puts $J_{\mathrm E}$ there.  This proves (G108.11).
If a prime of $B$ contains $J\cap J_{\mathrm E}+\lambda_\theta$,
primality makes it contain $J$ or $J_{\mathrm E}$, because their
product is contained in their intersection.  The corresponding
finite-weight condition makes that prime $\mathfrak m_B$.
The intersection ideal and $\lambda_\theta$ are contained in
$\mathfrak m_B$.  Hence the radical is exactly $\mathfrak m_B$,
proving the last finite-weight assertion.
\end{proof}

\paragraph{Exact sources and scope.}
The source commit is
\begin{center}
\texttt{7cbf7a0c32fca1badf05931879a37c3f8a2b03b4}.
\end{center}
Under \path{Development/AutomorphicArithmetic/AutomorphicVarieties/},
complete containing proofs read for this argument are
\path{CurveDVRSpecialization/FlipTwoThree.lean},
\path{CurveDVRSpecialization/NicePrimeEx.lean},
\path{CurveDVRSpecialization/LocProp.lean},
\path{CurveDVRSpecialization/GlobProp.lean},
\path{CurveDVRSpecialization/ChartOfFlat.lean},
\path{CurveDVRSpecialization/GlobOfCharts.lean},
\path{CurveDVRSpecialization/ExistsNicePrimeAboveOfLargeDim.lean},
and \path{CurveDVRSpecialization/PropagationTheoremGL3.lean}.
The graded specialization and local dimension interfaces are
\path{FrobeniusSmoothVectors/MinHom.lean},
\path{FrobeniusSmoothVectors/MonMemPow.lean},
\path{FrobeniusSmoothVectors/FibreClosed.lean},
\path{FrobeniusSmoothVectors/FrobeniusCornerLiftTwoToThree.lean},
\path{FrobeniusSmoothVectors/Seedpt.lean} and
\path{FrobeniusSmoothVectors/LocDim.lean}.
The proofs above establish the changed-tuple statements, including
the stronger normalized exclusion; they do not invoke the original
rank-three propagation theorem with an incompatible residual frame.
No Lean build or fresh audit of every foundational imported theorem
is asserted.  The usual dimension formula, flat completion and
Jacobson results remain the explicit commutative-algebra inputs
in the applied generic source lemmas.

Geometrically, the grading torus gives homogeneous component
specializations to the punctured closed fiber.  Finite homogeneous
pieces make the specialization survive Nakayama's lemma.  The
normalized collision subschemes are the actual bad loci removed
before choosing one-dimensional characteristic-$p$ specializations.
The coherent support $V(JA)$ propagates first through formal
component intersections and then over the connected core.  This is
an actual finite-type/formal-scheme argument; it constructs no
arithmetic surface square and makes no new Grothendieck archive
attribution.  Even after (E) and (S) are supplied, an application
to the original counting argument must verify its own character,
curve, local-length and omitted-prime inputs.  H04.4 and cubic RH
are not conclusions of this section.
